% ---------------------------------------------------------------------
\section{本章小结：对应表与常见误解}\label{sec:disall-summary}

本章的推理链是：非物理区的领头区域 $\to$ 嵌套的减除项 $\to$ 与 $\MSbar$ 对接 $\to$ 物理区（矩与固定 $x$ 两条途径）$\to$ \cref{thm:disall-factorization}（[A]）；与附录~N 的对应见\cref{tab:disall-correspondence}。

\begin{table}[htbp]
\centering
\caption{本章的陈述与附录~N 中精确形式的对应。“误差类型”按\cref{def:intro-error-types}：E-S 为标度陈述，E-N 为积分范数界，E-N$^\varphi$ 为纤维化积分范数界。}
\label{tab:disall-correspondence}
\small
\begin{tabularx}{\textwidth}{@{}>{\raggedright\arraybackslash}p{0.2\textwidth}>{\raggedright\arraybackslash}p{0.26\textwidth}>{\raggedright\arraybackslash}X@{}}
\toprule
本章陈述 & 附录~N 中的精确形式 & 差别（所用假设；误差类型） \\
\midrule
\cref{prop:disall-two-routes}（$R_{\mathrm{cond}}$） & \cref{sec:appN-processes}；\cref{thm:appN-subtraction}、\cref{lem:appN-error-table} &
(a)：\cref{asm:appN-collins}、\cref{asm:appN-contour}、\cref{asm:appN-envelope}、\cref{asm:appN-separation}；(b) 另需复圆盘上的一致性；(c) 另需\cref{asm:appN-cutsum}。E-N。 \\
\cref{prop:disall-factorized-subtraction}（$R_{\mathrm{cond}}$） & \cref{sec:appN-processes}；\cref{prop:method-product}、\cref{lem:appN-algebra} &
\cref{asm:appN-ward-product}(WI1)(WI2)、\cref{asm:appN-uv}；被积函数层次的恒等式（不含估计）。 \\
\cref{thm:disall-renorm-compat}（$R_{\mathrm{cond}}$） & \cref{sec:appN-processes} &
\cref{asm:appN-multren}(Z1)(Z3)(i)、\cref{asm:appN-ir-regulator}(Z2)、\cref{asm:appN-contour}(b)；内容等级 [C]，不被任何 [A] 陈述使用。 \\
\cref{ex:disall-ladder} & \cref{prop:appN-ladder}（\cref{sec:appN-examples}） &
代数精确（$R_{\mathrm{PT}}$）；费米子两级梯形图的 (CS1)--(CS3) 未证明；标量 $\varphi^3_6$ 模型单梯级梯形图的积分范数界（实 $x\ge x_0>1$）见\cref{prop:appN-phi36-dis}（E-N$^\varphi$）。 \\
\cref{lem:disall-fsi}、\cref{prop:disall-regions}（[A]） & \cref{asm:appN-cutsum}；\cref{asm:appN-envelope}、\cref{asm:appN-separation} &
标度陈述（E-S）及其割图形式与积分形式。 \\
\cref{thm:disall-factorization}（[A]） & \cref{thm:appN-structure}（$R_{\mathrm{cond}}$） &
条件形式列出附录~N 的全部技术假设，包括具名缺口\cref{asm:appN-brst-nonlocal}，故其内容等级为 [C]；本章的 [A] 定理只点名该缺口。文献：E-S；条件形式：E-N。 \\
\bottomrule
\end{tabularx}
\end{table}

% misconception box: M4 + critic O7; M25; M26
\begin{tcolorbox}[enhanced,breakable,title=本章涉及的常见误解,fonttitle=\bfseries,
  colback=gray!3!white,colframe=gray!55!black,boxrule=0.6pt,arc=2pt,
  left=5pt,right=5pt,top=3pt,bottom=3pt,before skip=6pt,after skip=6pt]
\begin{itemize}[leftmargin=1.2em,itemsep=1pt,topsep=1pt]
\item 「矩的因子化蕴含 $x$ 空间中一致的 $\order{\Lambda^2/Q^2}$ 余项」——矩只决定分布意义下的扭度 2 部分（\cref{cor:disall-consistency}）：对每个 $n$ 矩为 $\order{1/Q^2}$ 的 $\sin(Q^2x)$ 逐点为 $\order{1}$；
高扭度系数随 $n$ 增长，余项在 $x\to1$ 处不一致，共振区只能在平滑后比较。固定 $x$ 的逐点陈述是 [A] 定理\cref{thm:disall-factorization}，需要割图分析（\cref{prop:disall-two-routes}(c)）。
\item 「动量空间的 [A] 定理与 OPE 的 [R] 结果一致，所以逐点余项也是严格的」——一致的只是作为分布的扭度 2 部分（$R_{\mathrm{OPE}}$）；逐点余项只由全阶论证控制。
\item 「把 $\mu_R,\mu_F$ 变动 $2$ 倍就给出理论误差」——这是约定；定理只保证剩余的标度依赖高一阶（\cref{prop:disall-scale}）。
\item 「Callan--Gross 关系 $F_L=0$ 是 QCD 的定理」——它是部分子模型（领头阶）的结果；$F_L=\order{\as}$，$C_L^q=2C_Fz$，$C_L^g=4T_Fz(1-z)$（\cref{prop:disall-nlo}，Altarelli--Martinelli 关系）。
Drell--Yan 中的 Lam--Tung 关系也不是定理：它没有 $\order{\as}$ 修正\cite{Lam:1980uc}，但在 $\order{\as^2}$（$q_T\ne0$ 的次领头阶）被轻微破坏\cite{Mirkes:1992hu}。
\end{itemize}
\end{tcolorbox}

